\chapter{Results}
\label{ch:results}

% Report findings only. Interpretation belongs in the discussion.

\section{Literature Review}
\label{sec:results-lit}
This section reports the search and screening of Section~\ref{sec:lit-review}
and then what the screened literature says about the review question. All
counts come from the search scripts and the screening log; the corpus is listed
in Appendix~\ref{app:lit-inventory}.
\todo{The screening decisions in 01\_Literature/\allowbreak Literature-Review/\allowbreak screening.csv
were prepared by an AI assistant. Check every decision before this section
states who screened, and obtain the 244 full texts still missing.}

\subsection{Search and Screening}

\paragraph{Identification.} The searches ran on 28 October 2026; the final runs
started between 14:46 and 15:10~UTC (Appendix~\ref{app:lit-strings}).
Table~\ref{tab:lit-sources} gives the hits and records of every source. The four
databases returned 2{,}583 hits over the strands, which were 2{,}421 records
once each database's own repeats were merged; IEEE Xplore contributed most of
them (1{,}240), and its S1 query alone returned 1{,}004 hits. The eight venues
were listed completely: 25{,}196 papers, from 372 at SOUPS to 10{,}718 at CHI.
The venue rules selected 85 of them. The German PubMed strand (S5) returned no
record. Merging the database and venue records removed 111 duplicates across
sources and left 2{,}395 records from the searches. The other routes added 585
records that no search had found: 483 from citation searching alone, 64 from the
Methods Digest alone and 38 from the thesis citations alone or from more than
one of these routes.
In total, 2{,}980 records were screened. For 105 of them no abstract was
available, and 371 were screened with an abstract looked up by DOI.

\begin{table}[tbp]
\centering
\begin{singlespace}
\caption[Sources of the literature search and their yield]{Sources of the
literature search and their yield. Hits are counted per strand and include
repeats; records are counted after merging repeats within the source.
\emph{Only here}: records no other source found. \emph{Passed}: records that
passed title-and-abstract screening (Table~\ref{tab:lit-screening}).
$^a$~Strands applied locally. $^b$~Titles only.}
\label{tab:lit-sources}
\footnotesize
\setlength{\tabcolsep}{3.5pt}
\begin{tabularx}{\textwidth}{@{}>{\raggedright\arraybackslash}p{2.2cm}l>{\raggedright\arraybackslash}p{1.9cm}>{\raggedright\arraybackslash}Xrrr@{}}
\toprule
\textbf{Source} & \textbf{Kind} & \textbf{Fields} & \textbf{Hits} & \textbf{Records} & \textbf{Only here} & \textbf{Passed} \\
\midrule
% Generated by scripts/literature-search/v002/report.py from the frozen search and the screening log; do not edit.
PubMed & database & title, abstract & S1~634, S2~322, S3~13, S4~21, S5~0 & 911 & 850 & 61 \\
IEEE Xplore & database & title, abstract & S1~1{,}004, S2~170, S3~67, S4~56 & 1{,}240 & 1{,}154 & 84 \\
arXiv & database & title, abstract$^a$ & S1~119, S2~35, S3~91, S4~33 & 253 & 173 & 69 \\
ACL Anthology & database & title, abstract$^a$ & S1~1, S2~0, S3~14, S4~3 & 17 & 8 & 2 \\
\midrule
PoPETs & venue & title, abstract$^a$ & 1{,}160 listed & 27 & 12 & 10 \\
USENIX Security & venue & title, abstract$^a$ & 2{,}900 listed & 20 & 16 & 3 \\
SOUPS & venue & title, abstract$^a$ & 372 listed & 9 & 9 & 2 \\
CCS & venue & title$^{a,b}$ & 3{,}061 listed & 3 & 3 & 0 \\
IEEE S\&P & venue & title$^{a,b}$ & 1{,}818 listed & 2 & 0 & 2 \\
NDSS & venue & title$^{a,b}$ & 1{,}361 listed & 1 & 1 & 0 \\
CHI & venue & title$^{a,b}$ & 10{,}718 listed & 17 & 7 & 11 \\
CSCW & venue & title, abstract$^a$ & 3{,}806 listed & 6 & 4 & 1 \\
\midrule
Citation searching & other & references, citing works & 529 DOIs & 525 & 483 & 90 \\
Thesis citations & other & -- & 43 keys & 43 & 12 & 16 \\
Methods Digest & other & -- & 116 papers & 116 & 64 & 37 \\
\bottomrule
% the generated rows end with \bottomrule
\end{tabularx}
\end{singlespace}
\end{table}

Figure~\ref{fig:lit-sources} shows how much each source contributed beyond the
others. Most records of the three large databases were found by that database
alone:
1{,}154 of IEEE Xplore's 1{,}240, 850 of PubMed's 911 and 173 of arXiv's 253.
The overlap is larger where it matters for screening: of the 47 studies that
form the review evidence, 36 are in the Methods Digest and 17 were also found
on arXiv. Citation searching found the most records that passed screening
(90), followed by IEEE Xplore (84), arXiv (69) and PubMed (61). The venue lists
added few records, because the databases and citation searching had already
found most of their papers: CHI contributed 11 records that passed, PoPETs 10,
USENIX Security, SOUPS and IEEE S\&P together 7, CSCW 1, and CCS and NDSS none.

\begin{figure}[tbp]
\centering
% Generated by scripts/literature-search/v002/report.py from the frozen search and the screening log; do not edit.
\begin{tikzpicture}
\begin{axis}[xbar stacked, width=0.8\textwidth, height=92mm, bar width=6pt, xmin=0, xmax=1389,
  symbolic y coords={{Methods Digest},{Thesis citation},{Citation searching},{CSCW},{CHI},{NDSS},{IEEE S\&P},{CCS},{SOUPS},{USENIX Security},{PoPETs},{ACL Anthology},{arXiv},{IEEE Xplore},{PubMed}}, ytick=data, y tick label style={font=\footnotesize},
  x tick label style={font=\footnotesize, /pgf/number format/1000 sep={\,}},
  xlabel={Records after removing duplicates within the source}, ylabel={Source}, label style={font=\footnotesize},
  axis x line*=bottom, axis y line*=left, xmajorgrids, grid style={muted!25}, enlarge y limits=0.04,
  legend style={font=\scriptsize, draw=none, at={(0.98,0.03)}, anchor=south east}, legend cell align=left]
\addplot[fill=primary!75, draw=primary] coordinates {(64,{Methods Digest}) (12,{Thesis citation}) (483,{Citation searching}) (4,{CSCW}) (7,{CHI}) (1,{NDSS}) (0,{IEEE S\&P}) (3,{CCS}) (9,{SOUPS}) (16,{USENIX Security}) (12,{PoPETs}) (8,{ACL Anthology}) (173,{arXiv}) (1154,{IEEE Xplore}) (850,{PubMed})};
\addplot[fill=primary!18, draw=primary!60] coordinates {(52,{Methods Digest}) (31,{Thesis citation}) (42,{Citation searching}) (2,{CSCW}) (10,{CHI}) (0,{NDSS}) (2,{IEEE S\&P}) (0,{CCS}) (0,{SOUPS}) (4,{USENIX Security}) (15,{PoPETs}) (9,{ACL Anthology}) (80,{arXiv}) (86,{IEEE Xplore}) (61,{PubMed})};
\legend{found by this source only, also found by another source}
\node[anchor=west, font=\scriptsize, text=accentdark] at (axis cs:911,{PubMed}) {\,61};
\node[anchor=west, font=\scriptsize, text=accentdark] at (axis cs:1240,{IEEE Xplore}) {\,84};
\node[anchor=west, font=\scriptsize, text=accentdark] at (axis cs:253,{arXiv}) {\,69};
\node[anchor=west, font=\scriptsize, text=accentdark] at (axis cs:17,{ACL Anthology}) {\,2};
\node[anchor=west, font=\scriptsize, text=accentdark] at (axis cs:27,{PoPETs}) {\,10};
\node[anchor=west, font=\scriptsize, text=accentdark] at (axis cs:20,{USENIX Security}) {\,3};
\node[anchor=west, font=\scriptsize, text=accentdark] at (axis cs:9,{SOUPS}) {\,2};
\node[anchor=west, font=\scriptsize, text=accentdark] at (axis cs:3,{CCS}) {\,0};
\node[anchor=west, font=\scriptsize, text=accentdark] at (axis cs:2,{IEEE S\&P}) {\,2};
\node[anchor=west, font=\scriptsize, text=accentdark] at (axis cs:1,{NDSS}) {\,0};
\node[anchor=west, font=\scriptsize, text=accentdark] at (axis cs:17,{CHI}) {\,11};
\node[anchor=west, font=\scriptsize, text=accentdark] at (axis cs:6,{CSCW}) {\,1};
\node[anchor=west, font=\scriptsize, text=accentdark] at (axis cs:525,{Citation searching}) {\,90};
\node[anchor=west, font=\scriptsize, text=accentdark] at (axis cs:43,{Thesis citation}) {\,16};
\node[anchor=west, font=\scriptsize, text=accentdark] at (axis cs:116,{Methods Digest}) {\,37};
\end{axis}
\end{tikzpicture}

\caption[Records contributed by each source]{Records contributed by each
source, after merging repeats within the source (n~=~2{,}980 records; a record
found by several sources counts for each). Dark: found by this source only;
light: also found by another source. Numbers at the bars: records that passed
title-and-abstract screening (n~=~292).}
\Description{Horizontal stacked bars for 15 sources. IEEE Xplore and PubMed
have the longest bars, mostly unique records; citation searching follows.
Venue bars are short.}
\label{fig:lit-sources}
\end{figure}

\paragraph{Screening.} Table~\ref{tab:lit-screening} and
Figure~\ref{fig:lit-exclusions} give the screening by rule, and
Figure~\ref{fig:lit-flow} the whole flow. Each excluded record counts once,
under the first rule that applied, so each rule's denominator is the number of
records left after the rules before it. R3 excluded the most records (1{,}243
of 2{,}502, 49.7\,\%), for example security systems for hospitals, health
records or the Internet of Things in general. R4 removed 585 of 1{,}259 (46.5\,\%) studies of
health apps or wearables without a privacy question, R5 276 of 674 (40.9\,\%)
privacy-preserving systems for wearables and health apps, and R6 106 of 398
(26.6\,\%) studies of privacy text that neither used language models nor
analysed the wording of risk. R1 removed 423 records (14.2\,\%) that were not
research articles, among them 30 chapters and front or back matter of one
edited volume and eight retracted papers or retraction notices; R2 removed 55
works published before 2014. After R1--R6, 292 records passed to full text.

\begin{table}[tbp]
\centering
\begin{singlespace}
\caption[Screening of the literature review by rule]{Screening of the
literature review by rule, in the order applied. Each excluded record has one
primary reason. \emph{Assessed}: records that reached the rule; \emph{\%}:
excluded as a share of assessed. Records awaiting their full text leave the
count before R7 without being excluded.}
\label{tab:lit-screening}
\small
\begin{tabularx}{\textwidth}{@{}l>{\raggedright\arraybackslash}Xrrrr@{}}
\toprule
\textbf{Rule} & \textbf{Excludes} & \textbf{Assessed} & \textbf{Excluded} & \textbf{Remaining} & \textbf{\%} \\
\midrule
% Generated by scripts/literature-search/v002/report.py from the frozen search and the screening log; do not edit.
R1 & not a research article & 2{,}980 & 423 & 2{,}557 & 14.2 \\
R2 & outside language or period & 2{,}557 & 55 & 2{,}502 & 2.2 \\
R3 & outside the domain & 2{,}502 & 1{,}243 & 1{,}259 & 49.7 \\
R4 & domain without privacy focus & 1{,}259 & 585 & 674 & 46.5 \\
R5 & technical system without harms or communication & 674 & 276 & 398 & 40.9 \\
R6 & privacy text outside LQ3 and LQ4 & 398 & 106 & 292 & 26.6 \\
\midrule
-- & awaiting full text (not an exclusion) & 292 & \multicolumn{1}{r}{(244)} & 48 & -- \\
R7 & full text not retrievable & 48 & 0 & 48 & 0.0 \\
R8 & fails C1 or C4 on full text & 48 & 1 & 47 & 2.1 \\
\bottomrule
% the generated rows end with \bottomrule
\end{tabularx}
\end{singlespace}
\end{table}

\begin{figure}[tbp]
\centering
% Generated by scripts/literature-search/v002/report.py from the frozen search and the screening log; do not edit.
\begin{tikzpicture}
\begin{axis}[xbar, width=0.78\textwidth, height=70mm, bar width=6.5pt, xmin=0, xmax=1866,
  symbolic y coords={{R8 fails C1 or C4 on full text},{R7 full text not retrievable},{R6 privacy text outside LQ3 and LQ4},{R5 technical system without harms or communication},{R4 domain without privacy focus},{R3 outside the domain},{R2 outside language or period},{R1 not a research article}}, ytick=data,
  y tick label style={font=\footnotesize}, x tick label style={font=\footnotesize, /pgf/number format/1000 sep={\,}},
  xlabel={Records excluded (primary reason only)}, ylabel={Exclusion rule, in the order applied}, label style={font=\footnotesize},
  axis x line*=bottom, axis y line*=left, xmajorgrids, grid style={muted!25}, enlarge y limits=0.062]
\addplot[fill=primary!70, draw=primary] coordinates {(1,{R8 fails C1 or C4 on full text}) (0,{R7 full text not retrievable}) (106,{R6 privacy text outside LQ3 and LQ4}) (276,{R5 technical system without harms or communication}) (585,{R4 domain without privacy focus}) (1243,{R3 outside the domain}) (55,{R2 outside language or period}) (423,{R1 not a research article})};
\node[anchor=west, font=\scriptsize, text=muted] at (axis cs:423,{R1 not a research article}) {\,423 of 2{,}980 (14.2\,\%)};
\node[anchor=west, font=\scriptsize, text=muted] at (axis cs:55,{R2 outside language or period}) {\,55 of 2{,}557 (2.2\,\%)};
\node[anchor=west, font=\scriptsize, text=muted] at (axis cs:1243,{R3 outside the domain}) {\,1{,}243 of 2{,}502 (49.7\,\%)};
\node[anchor=west, font=\scriptsize, text=muted] at (axis cs:585,{R4 domain without privacy focus}) {\,585 of 1{,}259 (46.5\,\%)};
\node[anchor=west, font=\scriptsize, text=muted] at (axis cs:276,{R5 technical system without harms or communication}) {\,276 of 674 (40.9\,\%)};
\node[anchor=west, font=\scriptsize, text=muted] at (axis cs:106,{R6 privacy text outside LQ3 and LQ4}) {\,106 of 398 (26.6\,\%)};
\node[anchor=west, font=\scriptsize, text=muted] at (axis cs:0,{R7 full text not retrievable}) {\,0 of 48 (0.0\,\%)};
\node[anchor=west, font=\scriptsize, text=muted] at (axis cs:1,{R8 fails C1 or C4 on full text}) {\,1 of 48 (2.1\,\%)};
\end{axis}
\end{tikzpicture}

\caption[Exclusions by rule]{Records excluded by each rule, with the number
the rule assessed and the share excluded (n~=~2{,}980 screened; 2{,}689
excluded in total). R1--R6: title and abstract; R7--R8: full text (48
assessed).}
\Description{Horizontal bars for rules R1 to R8. R3 is the longest bar
(1,243), followed by R4 (585), R1 (423), R5 (276), R6 (106), R2 (55), R8 (1)
and R7 (0).}
\label{fig:lit-exclusions}
\end{figure}

\begin{figure}[tbp]
\centering
\resizebox{\textwidth}{!}{% Generated by scripts/literature-search/v002/report.py from the frozen search and the screening log; do not edit.
\begin{tikzpicture}[font=\footnotesize, node distance=5mm and 7mm,
  box/.style={draw=primary, fill=primary!5, rounded corners=1.5pt, align=left, text width=62mm, inner sep=3pt},
  side/.style={draw=muted, fill=white, rounded corners=1.5pt, align=left, text width=52mm, inner sep=3pt},
  stage/.style={font=\scriptsize\bfseries, text=primary, rotate=90, anchor=south},
  arr/.style={-{Stealth[length=1.5mm]}, semithick, accentdark}]
\node[box] (db) {\textbf{Database searches:} 2{,}421 records\\
  PubMed 911, IEEE Xplore 1{,}240,
  arXiv 253, ACL Anthology 17\\[2pt]
  \textbf{Venue searches:} 85 records\\
  PoPETs 27, USENIX Security 20, SOUPS 9, CCS 3, IEEE S\&P 2, NDSS 1, CHI 17, CSCW 6};
\node[side, right=of db] (oth) {\textbf{Other methods:} 585 records found\\ by no search\\
  citation searching 489, thesis\\ citations 33, Methods Digest 90$^{*}$};
\node[box, below=of db] (dd) {Duplicates across sources removed: 111\\
  Records from searches: 2{,}395};
\node[box, below=8mm of dd] (ta) {\textbf{Screened on title and abstract:} 2{,}980};
\node[side, right=of ta] (tax) {Excluded (R1--R6): 2{,}688\\
  R1 423, R2 55, R3 1{,}243,\\
  R4 585, R5 276, R6 106};
\node[box, below=of ta] (pass) {Passed to full text: 292};
\node[side, right=of pass] (pend) {Awaiting full text: 244\\
  LQ1 154, LQ2 67, LQ3 37, LQ4 21$^{\dagger}$};
\node[box, below=of pass] (ftb) {\textbf{Full texts assessed:} 48};
\node[side, right=of ftb] (ftx) {Excluded (R7--R8): 1\\ R7 0, R8 1};
\node[box, below=of ftb, fill=accent!10, draw=accent] (inc) {\textbf{Review evidence:} 47 studies\\
  LQ1 19, LQ2 21, LQ3 14, LQ4 6$^{\dagger}$};
\node[side, right=of inc, draw=accent] (ctx) {\textbf{Contextual collection:} 32\\ cited works excluded by R1--R8};
\draw[arr] (db) -- (dd);
\draw[arr] (dd) -- (ta);
\draw[arr] (oth.south) |- ([yshift=4mm]ta.north) -- (ta.north);
\draw[arr] (ta) -- (tax); \draw[arr] (ta) -- (pass); \draw[arr] (pass) -- (pend); \draw[arr] (pass) -- (ftb);
\draw[arr] (ftb) -- (ftx); \draw[arr] (ftb) -- (inc);
\draw[arr, dashed, muted] (tax.east) -- ++(3mm,0) |- (ctx.east);
\node[stage] at ([xshift=-2mm]db.west) {Identification};
\node[stage] at ([xshift=-2mm]dd.west) {Duplicates};
\node[stage] at ([xshift=-2mm, yshift=-4mm]pass.west) {Screening};
\node[stage] at ([xshift=-2mm]inc.west) {Included};
\end{tikzpicture}
}
\caption[Flow of records through the literature review]{Flow of records
through the literature review, after PRISMA~2020~\parencite{page2021prisma}.
$^{*}$~Records found by no search; several routes can share a record.
$^{\dagger}$~A study can answer several sub-questions. The dashed arrow marks
the 32 cited works that the rules excluded; they form the contextual
collection.}
\Description{Flow diagram from identification (database searches 2,421 and venue
searches 85 records) through the removal of 111 duplicates to 2,980 screened
records, 2,688 exclusions at title and abstract, 292 records passed to full
text, of which 244 await their full text and 48 were assessed; 47 studies are
review evidence. A side box lists the contextual collection of 32 cited
works.}
\label{fig:lit-flow}
\end{figure}

\paragraph{Full texts.} Full texts were assessed for the 48 records whose PDF
the project already holds. One was excluded under R8: it discusses the use of
language models for assessing policies without evaluating
them~\parencite{aydin2024assessing}. The other 47 are the review
evidence: 19 answer LQ1, 21 LQ2, 14 LQ3 and 6 LQ4. The full texts of the other
244 records have not been obtained yet; they are open access in most cases
(arXiv, PubMed Central, USENIX, PoPETs) but were not downloaded for this draft.
They remain in the corpus as \emph{awaiting full text}, with the sub-questions
assigned at title and abstract (154 LQ1, 67 LQ2, 37 LQ3, 21 LQ4), and no
statement below rests on them alone.

\paragraph{Existing literature.} All 43 works cited in the thesis before this
review entered the corpus. Sixteen meet the criteria and are review evidence.
The other 27 are in the contextual collection: 10 are not research articles
(R1; web pages, official statistics, a book, a book chapter and software), 5
predate 2014 (R2; among them Solove's taxonomy and Calo's boundaries of privacy
harm), 7 lie outside the domain (R3; among them Citron and Solove's typology,
the harm model of Knabenhans et al.\ and the work on annotation with language
models), and 5 analyse privacy text without language models or risk wording
(R6; OPP-115, Polisis, MAPS, the Data Safety study of Arkalakis et al.\ and the
analysis of Apple's privacy labels by Li et~al.). Five works cited since join
them: the four reporting and review guidelines of
Section~\ref{sec:lit-review} and the paper excluded at full text. Of the 116 Methods Digest
papers, 36 are review evidence, 19 are contextual (cited in the thesis), 60
were excluded and 1 awaits its full text. Citation searching from the four seed
papers yielded 529 DOIs, of which 4 had no metadata; 81 of its 525 records
passed screening, 61 of them found by no other route.

\subsection{What the Literature Shows}

\paragraph{Years and venues.} Figure~\ref{fig:lit-years} shows that the field
is young: of the 291 records that passed screening and were not excluded at full
text, 146 (50\,\%) were published from 2023 onwards, and although S3 searched
from 2022, the first LQ3 records date from 2023. The records come from a wide
range of venues
(Figure~\ref{fig:lit-venues}): 57 from health or medical journals, 51 from
security and privacy venues, 28 from HCI venues and 40 still only as
preprints. Of the 47 studies of the review evidence, 13 appeared in health or
medical journals, 11 at security and privacy venues and 6 at HCI venues; 10
are preprints.

\begin{figure}[tbp]
\centering
% Generated by scripts/literature-search/v002/report.py from the frozen search and the screening log; do not edit.
\begin{tikzpicture}
\begin{axis}[ybar stacked, width=0.92\textwidth, height=64mm, bar width=9pt, ymin=0,
  xtick={2014,2015,2016,2017,2018,2019,2020,2021,2022,2023,2024,2025,2026}, x tick label style={font=\footnotesize, rotate=45, anchor=east, /pgf/number format/1000 sep={}},
  y tick label style={font=\footnotesize}, xlabel={Year of publication}, ylabel={Records}, label style={font=\footnotesize},
  axis x line*=bottom, axis y line*=left, ymajorgrids, grid style={muted!25}, enlarge x limits=0.05,
  legend style={font=\scriptsize, draw=none, at={(0.02,0.98)}, anchor=north west}, legend cell align=left]
\addplot[fill=primary!80, draw=black!35] coordinates {(2014,2) (2015,5) (2016,15) (2017,11) (2018,18) (2019,11) (2020,17) (2021,14) (2022,15) (2023,14) (2024,23) (2025,17) (2026,11)};
\addplot[fill=accent!75, draw=black!35] coordinates {(2014,1) (2015,1) (2016,0) (2017,1) (2018,4) (2019,4) (2020,6) (2021,5) (2022,2) (2023,6) (2024,2) (2025,9) (2026,9)};
\addplot[fill=linkblue!50, draw=black!35] coordinates {(2014,0) (2015,0) (2016,0) (2017,0) (2018,0) (2019,0) (2020,0) (2021,0) (2022,0) (2023,3) (2024,14) (2025,19) (2026,14)};
\addplot[fill=muted!40, draw=black!35] coordinates {(2014,1) (2015,0) (2016,2) (2017,1) (2018,2) (2019,1) (2020,2) (2021,4) (2022,0) (2023,1) (2024,3) (2025,0) (2026,1)};
\legend{LQ1 harms, LQ2 communication, LQ3 language models, LQ4 framing}
\end{axis}
\end{tikzpicture}

\par\smallskip
{\scriptsize\setlength{\tabcolsep}{3.1pt}
\begin{tabular}{@{}l*{13}{r}@{}}
Year & 2014 & 2015 & 2016 & 2017 & 2018 & 2019 & 2020 & 2021 & 2022 & 2023 & 2024 & 2025 & 2026 \\
Review evidence & % Generated by scripts/literature-search/v002/report.py from the frozen search and the screening log; do not edit.
0 & 3 & 0 & 0 & 3 & 6 & 3 & 2 & 3 & 6 & 8 & 7 & 6 \\

\end{tabular}}
\caption[Publication years of the screened literature]{Publication years of the
records that passed title-and-abstract screening and were not excluded at full
text (n~=~291: 47 review evidence and 244 awaiting full text), by first
sub-question. The row below gives the 47 studies of the review evidence. 2026
covers January to 28 October.}
\Description{Stacked bars per year from 2014 to 2026. Counts rise from 4 in 2014
to 42 in 2024 and 45 in 2025; LQ1 dominates every year, LQ3 appears from 2023.}
\label{fig:lit-years}
\end{figure}

\begin{figure}[tbp]
\centering
% Generated by scripts/literature-search/v002/report.py from the frozen search and the screening log; do not edit.
\begin{tikzpicture}
\begin{axis}[xbar stacked, width=0.75\textwidth, height=76mm, bar width=7pt, xmin=0,
  xmax=112, symbolic y coords={{Preprint},{Other computing venue},{Law or social science},{NLP venue},{Health or medical journal},{HCI venue},{Security and privacy venue}}, ytick=data,
  y tick label style={font=\footnotesize}, x tick label style={font=\footnotesize}, xlabel={Records}, ylabel={Venue type},
  label style={font=\footnotesize}, axis x line*=bottom, axis y line*=left, xmajorgrids, grid style={muted!25},
  enlarge y limits=0.071, legend style={font=\scriptsize, draw=none, at={(0.98,0.03)}, anchor=south east}, legend cell align=left]
\addplot[fill=accent!80, draw=accent] coordinates {(10,{Preprint}) (5,{Other computing venue}) (2,{Law or social science}) (0,{NLP venue}) (13,{Health or medical journal}) (6,{HCI venue}) (11,{Security and privacy venue})};
\addplot[fill=accent!20, draw=accent!60] coordinates {(30,{Preprint}) (92,{Other computing venue}) (15,{Law or social science}) (1,{NLP venue}) (44,{Health or medical journal}) (22,{HCI venue}) (40,{Security and privacy venue})};
\legend{review evidence, awaiting full text}
\end{axis}
\end{tikzpicture}

\caption[Venue types of the screened literature]{Venue types of the records that
passed screening (n~=~291), coded from the venue name with fixed keyword
rules (\texttt{extraction\_rules} in the search configuration). Dark: review
evidence (n~=~47); light: awaiting full text (n~=~244).}
\Description{Horizontal stacked bars for seven venue types; other computing
venues, health journals and security venues are the largest groups.}
\label{fig:lit-venues}
\end{figure}

\paragraph{Evidence map.} Figure~\ref{fig:lit-map-surfaces} crosses the app
domain with the material each study analysed, coded from titles and abstracts
with fixed keyword rules for all 291 records. Wearables in general (90) and
activity and fitness (85) dominate; reproductive health follows with 44,
mental health with 23, and sleep and diet with 9 and 10. Privacy policies are
the most analysed surface (139 records), followed by traffic or code (68) and
interfaces (54). Store listings appear in two records only, one on mental
health apps and one on medical apps, and privacy labels in ten. Of the 88
records answering LQ2, 43 analyse one surface only. The communication of health
tracking apps is thus studied mostly through one document, the policy, and its
relation to app behaviour; the store listing, which every user sees before
installing, is almost unstudied.

\begin{figure}[tbp]
\centering
% Generated by scripts/literature-search/v002/report.py from the frozen search and the screening log; do not edit.
\begin{tikzpicture}[font=\footnotesize, x=9.5mm, y=-6.2mm]
\node[rotate=45, anchor=south west, inner sep=1pt, font=\scriptsize] at (0.3,-0.05) {Privacy policy};
\node[rotate=45, anchor=south west, inner sep=1pt, font=\scriptsize] at (1.3,-0.05) {Privacy label};
\node[rotate=45, anchor=south west, inner sep=1pt, font=\scriptsize] at (2.3,-0.05) {Store listing};
\node[rotate=45, anchor=south west, inner sep=1pt, font=\scriptsize] at (3.3,-0.05) {App interface};
\node[rotate=45, anchor=south west, inner sep=1pt, font=\scriptsize] at (4.3,-0.05) {Marketing or website};
\node[rotate=45, anchor=south west, inner sep=1pt, font=\scriptsize] at (5.3,-0.05) {Traffic or code};
\node[rotate=45, anchor=south west, inner sep=1pt, font=\scriptsize] at (6.3,-0.05) {Sensor data};
\node[rotate=45, anchor=south west, inner sep=1pt, font=\scriptsize] at (7.3,-0.05) {Users' accounts};
\node[rotate=45, anchor=south west, inner sep=1pt, font=\scriptsize] at (8.3,-0.05) {Other notice};
\node[rotate=45, anchor=south west, inner sep=1pt, font=\scriptsize] at (9.3,-0.05) {Not stated};
\node[anchor=east, font=\scriptsize] at (-0.1,0.5) {Activity and fitness};
\filldraw[fill=primary!29, draw=muted!35] (0,0) rectangle (1,1);
\node[text=black, font=\scriptsize] at (0.5,0.5) {15};
\filldraw[fill=primary!11, draw=muted!35] (1,0) rectangle (2,1);
\node[text=black, font=\scriptsize] at (1.5,0.5) {1};
\filldraw[fill=primary!11, draw=muted!35] (2,0) rectangle (3,1);
\node[text=black, font=\scriptsize] at (2.5,0.5) {1};
\filldraw[fill=primary!35, draw=muted!35] (3,0) rectangle (4,1);
\node[text=black, font=\scriptsize] at (3.5,0.5) {20};
\filldraw[fill=primary!18, draw=muted!35] (4,0) rectangle (5,1);
\node[text=black, font=\scriptsize] at (4.5,0.5) {7};
\filldraw[fill=primary!48, draw=muted!35] (5,0) rectangle (6,1);
\node[text=black, font=\scriptsize] at (5.5,0.5) {30};
\filldraw[fill=primary!37, draw=muted!35] (6,0) rectangle (7,1);
\node[text=black, font=\scriptsize] at (6.5,0.5) {22};
\filldraw[fill=primary!32, draw=muted!35] (7,0) rectangle (8,1);
\node[text=black, font=\scriptsize] at (7.5,0.5) {18};
\filldraw[fill=primary!12, draw=muted!35] (8,0) rectangle (9,1);
\node[text=black, font=\scriptsize] at (8.5,0.5) {2};
\filldraw[fill=primary!31, draw=muted!35] (9,0) rectangle (10,1);
\node[text=black, font=\scriptsize] at (9.5,0.5) {17};
\node[anchor=east, font=\scriptsize] at (-0.1,1.5) {Reproductive health};
\filldraw[fill=primary!40, draw=muted!35] (0,1) rectangle (1,2);
\node[text=black, font=\scriptsize] at (0.5,1.5) {24};
\filldraw[fill=primary!12, draw=muted!35] (1,1) rectangle (2,2);
\node[text=black, font=\scriptsize] at (1.5,1.5) {2};
\filldraw[fill=white, draw=muted!35] (2,1) rectangle (3,2);
\filldraw[fill=primary!20, draw=muted!35] (3,1) rectangle (4,2);
\node[text=black, font=\scriptsize] at (3.5,1.5) {8};
\filldraw[fill=primary!17, draw=muted!35] (4,1) rectangle (5,2);
\node[text=black, font=\scriptsize] at (4.5,1.5) {6};
\filldraw[fill=primary!23, draw=muted!35] (5,1) rectangle (6,2);
\node[text=black, font=\scriptsize] at (5.5,1.5) {11};
\filldraw[fill=white, draw=muted!35] (6,1) rectangle (7,2);
\filldraw[fill=primary!22, draw=muted!35] (7,1) rectangle (8,2);
\node[text=black, font=\scriptsize] at (7.5,1.5) {10};
\filldraw[fill=primary!15, draw=muted!35] (8,1) rectangle (9,2);
\node[text=black, font=\scriptsize] at (8.5,1.5) {4};
\filldraw[fill=primary!22, draw=muted!35] (9,1) rectangle (10,2);
\node[text=black, font=\scriptsize] at (9.5,1.5) {10};
\node[anchor=east, font=\scriptsize] at (-0.1,2.5) {Mental health};
\filldraw[fill=primary!29, draw=muted!35] (0,2) rectangle (1,3);
\node[text=black, font=\scriptsize] at (0.5,2.5) {15};
\filldraw[fill=primary!11, draw=muted!35] (1,2) rectangle (2,3);
\node[text=black, font=\scriptsize] at (1.5,2.5) {1};
\filldraw[fill=primary!12, draw=muted!35] (2,2) rectangle (3,3);
\node[text=black, font=\scriptsize] at (2.5,2.5) {2};
\filldraw[fill=primary!18, draw=muted!35] (3,2) rectangle (4,3);
\node[text=black, font=\scriptsize] at (3.5,2.5) {7};
\filldraw[fill=primary!15, draw=muted!35] (4,2) rectangle (5,3);
\node[text=black, font=\scriptsize] at (4.5,2.5) {4};
\filldraw[fill=primary!20, draw=muted!35] (5,2) rectangle (6,3);
\node[text=black, font=\scriptsize] at (5.5,2.5) {8};
\filldraw[fill=primary!11, draw=muted!35] (6,2) rectangle (7,3);
\node[text=black, font=\scriptsize] at (6.5,2.5) {1};
\filldraw[fill=primary!13, draw=muted!35] (7,2) rectangle (8,3);
\node[text=black, font=\scriptsize] at (7.5,2.5) {3};
\filldraw[fill=primary!11, draw=muted!35] (8,2) rectangle (9,3);
\node[text=black, font=\scriptsize] at (8.5,2.5) {1};
\filldraw[fill=primary!15, draw=muted!35] (9,2) rectangle (10,3);
\node[text=black, font=\scriptsize] at (9.5,2.5) {4};
\node[anchor=east, font=\scriptsize] at (-0.1,3.5) {Sleep};
\filldraw[fill=white, draw=muted!35] (0,3) rectangle (1,4);
\filldraw[fill=white, draw=muted!35] (1,3) rectangle (2,4);
\filldraw[fill=white, draw=muted!35] (2,3) rectangle (3,4);
\filldraw[fill=primary!12, draw=muted!35] (3,3) rectangle (4,4);
\node[text=black, font=\scriptsize] at (3.5,3.5) {2};
\filldraw[fill=white, draw=muted!35] (4,3) rectangle (5,4);
\filldraw[fill=primary!13, draw=muted!35] (5,3) rectangle (6,4);
\node[text=black, font=\scriptsize] at (5.5,3.5) {3};
\filldraw[fill=white, draw=muted!35] (6,3) rectangle (7,4);
\filldraw[fill=primary!11, draw=muted!35] (7,3) rectangle (8,4);
\node[text=black, font=\scriptsize] at (7.5,3.5) {1};
\filldraw[fill=white, draw=muted!35] (8,3) rectangle (9,4);
\filldraw[fill=primary!15, draw=muted!35] (9,3) rectangle (10,4);
\node[text=black, font=\scriptsize] at (9.5,3.5) {4};
\node[anchor=east, font=\scriptsize] at (-0.1,4.5) {Diet and weight};
\filldraw[fill=primary!16, draw=muted!35] (0,4) rectangle (1,5);
\node[text=black, font=\scriptsize] at (0.5,4.5) {5};
\filldraw[fill=white, draw=muted!35] (1,4) rectangle (2,5);
\filldraw[fill=white, draw=muted!35] (2,4) rectangle (3,5);
\filldraw[fill=primary!12, draw=muted!35] (3,4) rectangle (4,5);
\node[text=black, font=\scriptsize] at (3.5,4.5) {2};
\filldraw[fill=white, draw=muted!35] (4,4) rectangle (5,5);
\filldraw[fill=primary!11, draw=muted!35] (5,4) rectangle (6,5);
\node[text=black, font=\scriptsize] at (5.5,4.5) {1};
\filldraw[fill=primary!12, draw=muted!35] (6,4) rectangle (7,5);
\node[text=black, font=\scriptsize] at (6.5,4.5) {2};
\filldraw[fill=primary!11, draw=muted!35] (7,4) rectangle (8,5);
\node[text=black, font=\scriptsize] at (7.5,4.5) {1};
\filldraw[fill=primary!11, draw=muted!35] (8,4) rectangle (9,5);
\node[text=black, font=\scriptsize] at (8.5,4.5) {1};
\filldraw[fill=primary!12, draw=muted!35] (9,4) rectangle (10,5);
\node[text=black, font=\scriptsize] at (9.5,4.5) {2};
\node[anchor=east, font=\scriptsize] at (-0.1,5.5) {Wearables in general};
\filldraw[fill=primary!18, draw=muted!35] (0,5) rectangle (1,6);
\node[text=black, font=\scriptsize] at (0.5,5.5) {7};
\filldraw[fill=white, draw=muted!35] (1,5) rectangle (2,6);
\filldraw[fill=white, draw=muted!35] (2,5) rectangle (3,6);
\filldraw[fill=primary!29, draw=muted!35] (3,5) rectangle (4,6);
\node[text=black, font=\scriptsize] at (3.5,5.5) {15};
\filldraw[fill=primary!13, draw=muted!35] (4,5) rectangle (5,6);
\node[text=black, font=\scriptsize] at (4.5,5.5) {3};
\filldraw[fill=primary!40, draw=muted!35] (5,5) rectangle (6,6);
\node[text=black, font=\scriptsize] at (5.5,5.5) {24};
\filldraw[fill=primary!60, draw=muted!35] (6,5) rectangle (7,6);
\node[text=white, font=\scriptsize] at (6.5,5.5) {40};
\filldraw[fill=primary!27, draw=muted!35] (7,5) rectangle (8,6);
\node[text=black, font=\scriptsize] at (7.5,5.5) {14};
\filldraw[fill=primary!12, draw=muted!35] (8,5) rectangle (9,6);
\node[text=black, font=\scriptsize] at (8.5,5.5) {2};
\filldraw[fill=primary!34, draw=muted!35] (9,5) rectangle (10,6);
\node[text=black, font=\scriptsize] at (9.5,5.5) {19};
\node[anchor=east, font=\scriptsize] at (-0.1,6.5) {Health apps in general};
\filldraw[fill=primary!37, draw=muted!35] (0,6) rectangle (1,7);
\node[text=black, font=\scriptsize] at (0.5,6.5) {22};
\filldraw[fill=primary!12, draw=muted!35] (1,6) rectangle (2,7);
\node[text=black, font=\scriptsize] at (1.5,6.5) {2};
\filldraw[fill=white, draw=muted!35] (2,6) rectangle (3,7);
\filldraw[fill=primary!20, draw=muted!35] (3,6) rectangle (4,7);
\node[text=black, font=\scriptsize] at (3.5,6.5) {8};
\filldraw[fill=primary!16, draw=muted!35] (4,6) rectangle (5,7);
\node[text=black, font=\scriptsize] at (4.5,6.5) {5};
\filldraw[fill=primary!23, draw=muted!35] (5,6) rectangle (6,7);
\node[text=black, font=\scriptsize] at (5.5,6.5) {11};
\filldraw[fill=primary!12, draw=muted!35] (6,6) rectangle (7,7);
\node[text=black, font=\scriptsize] at (6.5,6.5) {2};
\filldraw[fill=primary!20, draw=muted!35] (7,6) rectangle (8,7);
\node[text=black, font=\scriptsize] at (7.5,6.5) {8};
\filldraw[fill=primary!12, draw=muted!35] (8,6) rectangle (9,7);
\node[text=black, font=\scriptsize] at (8.5,6.5) {2};
\filldraw[fill=primary!13, draw=muted!35] (9,6) rectangle (10,7);
\node[text=black, font=\scriptsize] at (9.5,6.5) {3};
\node[anchor=east, font=\scriptsize] at (-0.1,7.5) {Outside health};
\filldraw[fill=primary!85, draw=muted!35] (0,7) rectangle (1,8);
\node[text=white, font=\scriptsize] at (0.5,7.5) {59};
\filldraw[fill=primary!16, draw=muted!35] (1,7) rectangle (2,8);
\node[text=black, font=\scriptsize] at (1.5,7.5) {5};
\filldraw[fill=white, draw=muted!35] (2,7) rectangle (3,8);
\filldraw[fill=primary!17, draw=muted!35] (3,7) rectangle (4,8);
\node[text=black, font=\scriptsize] at (3.5,7.5) {6};
\filldraw[fill=primary!36, draw=muted!35] (4,7) rectangle (5,8);
\node[text=black, font=\scriptsize] at (4.5,7.5) {21};
\filldraw[fill=primary!20, draw=muted!35] (5,7) rectangle (6,8);
\node[text=black, font=\scriptsize] at (5.5,7.5) {8};
\filldraw[fill=primary!11, draw=muted!35] (6,7) rectangle (7,8);
\node[text=black, font=\scriptsize] at (6.5,7.5) {1};
\filldraw[fill=primary!17, draw=muted!35] (7,7) rectangle (8,8);
\node[text=black, font=\scriptsize] at (7.5,7.5) {6};
\filldraw[fill=primary!22, draw=muted!35] (8,7) rectangle (9,8);
\node[text=black, font=\scriptsize] at (8.5,7.5) {10};
\filldraw[fill=primary!12, draw=muted!35] (9,7) rectangle (10,8);
\node[text=black, font=\scriptsize] at (9.5,7.5) {2};
\node[anchor=north, font=\footnotesize] at (5.0,8.2) {Material or surface studied (title and abstract)};
\node[anchor=south, rotate=90, font=\footnotesize] at (-3.58,4.0) {App domain};
\end{tikzpicture}

\caption[Evidence map: app domain and material studied]{Evidence map of the
records that passed screening (n~=~291): number of records per app domain and
material or surface, coded from title and abstract with fixed keyword rules;
a record can fall in several cells. \emph{Outside health}: LQ3 and LQ4 studies
of privacy text in other domains.}
\Description{Heat map with eight app domains as rows and ten surfaces as
columns. The darkest cells are privacy policies in studies outside health,
sensor data of wearables and traffic of activity apps; the store-listing column
is almost empty.}
\label{fig:lit-map-surfaces}
\end{figure}

Figure~\ref{fig:lit-map-rq} shows, for the 47 studies of the review evidence,
which part of the thesis RQ each bears on, coded by hand from the full text.
Twenty-six name or test harms (RQ1), ten examine protection claims (RQ3) and
eleven compare several surfaces (RQ4); only five address how harms are worded
or downplayed (RQ2), and of these only the work on privacy
washing~\parencite{brackin2026privacy} also examines protection claims. No study of the review evidence
addresses all four parts, and none codes stance or the framing of PETs in
the communication of health tracking apps. Thirty-five serve as method
precedents.

\begin{figure}[tbp]
\centering
% Generated by scripts/literature-search/v002/report.py from the frozen search and the screening log; do not edit.
\begin{tikzpicture}[font=\footnotesize, x=9.5mm, y=-6.2mm]
\node[rotate=45, anchor=south west, inner sep=1pt, font=\scriptsize] at (0.3,-0.05) {Activity and fitness};
\node[rotate=45, anchor=south west, inner sep=1pt, font=\scriptsize] at (1.3,-0.05) {Reproductive health};
\node[rotate=45, anchor=south west, inner sep=1pt, font=\scriptsize] at (2.3,-0.05) {Mental health};
\node[rotate=45, anchor=south west, inner sep=1pt, font=\scriptsize] at (3.3,-0.05) {Sleep};
\node[rotate=45, anchor=south west, inner sep=1pt, font=\scriptsize] at (4.3,-0.05) {Diet and weight};
\node[rotate=45, anchor=south west, inner sep=1pt, font=\scriptsize] at (5.3,-0.05) {Wearables in general};
\node[rotate=45, anchor=south west, inner sep=1pt, font=\scriptsize] at (6.3,-0.05) {Health apps in general};
\node[rotate=45, anchor=south west, inner sep=1pt, font=\scriptsize] at (7.3,-0.05) {Outside health};
\node[anchor=east, font=\scriptsize] at (-0.1,0.5) {RQ1 harms named or missing};
\filldraw[fill=primary!43, draw=muted!35] (0,0) rectangle (1,1);
\node[text=black, font=\scriptsize] at (0.5,0.5) {8};
\filldraw[fill=primary!39, draw=muted!35] (1,0) rectangle (2,1);
\node[text=black, font=\scriptsize] at (1.5,0.5) {7};
\filldraw[fill=primary!39, draw=muted!35] (2,0) rectangle (3,1);
\node[text=black, font=\scriptsize] at (2.5,0.5) {7};
\filldraw[fill=primary!14, draw=muted!35] (3,0) rectangle (4,1);
\node[text=black, font=\scriptsize] at (3.5,0.5) {1};
\filldraw[fill=white, draw=muted!35] (4,0) rectangle (5,1);
\filldraw[fill=primary!22, draw=muted!35] (5,0) rectangle (6,1);
\node[text=black, font=\scriptsize] at (5.5,0.5) {3};
\filldraw[fill=primary!18, draw=muted!35] (6,0) rectangle (7,1);
\node[text=black, font=\scriptsize] at (6.5,0.5) {2};
\filldraw[fill=primary!22, draw=muted!35] (7,0) rectangle (8,1);
\node[text=black, font=\scriptsize] at (7.5,0.5) {3};
\node[anchor=east, font=\scriptsize] at (-0.1,1.5) {RQ2 stance on harms};
\filldraw[fill=white, draw=muted!35] (0,1) rectangle (1,2);
\filldraw[fill=primary!14, draw=muted!35] (1,1) rectangle (2,2);
\node[text=black, font=\scriptsize] at (1.5,1.5) {1};
\filldraw[fill=white, draw=muted!35] (2,1) rectangle (3,2);
\filldraw[fill=white, draw=muted!35] (3,1) rectangle (4,2);
\filldraw[fill=white, draw=muted!35] (4,1) rectangle (5,2);
\filldraw[fill=white, draw=muted!35] (5,1) rectangle (6,2);
\filldraw[fill=white, draw=muted!35] (6,1) rectangle (7,2);
\filldraw[fill=primary!26, draw=muted!35] (7,1) rectangle (8,2);
\node[text=black, font=\scriptsize] at (7.5,1.5) {4};
\node[anchor=east, font=\scriptsize] at (-0.1,2.5) {RQ3 protection claims};
\filldraw[fill=primary!22, draw=muted!35] (0,2) rectangle (1,3);
\node[text=black, font=\scriptsize] at (0.5,2.5) {3};
\filldraw[fill=primary!22, draw=muted!35] (1,2) rectangle (2,3);
\node[text=black, font=\scriptsize] at (1.5,2.5) {3};
\filldraw[fill=primary!18, draw=muted!35] (2,2) rectangle (3,3);
\node[text=black, font=\scriptsize] at (2.5,2.5) {2};
\filldraw[fill=white, draw=muted!35] (3,2) rectangle (4,3);
\filldraw[fill=white, draw=muted!35] (4,2) rectangle (5,3);
\filldraw[fill=primary!14, draw=muted!35] (5,2) rectangle (6,3);
\node[text=black, font=\scriptsize] at (5.5,2.5) {1};
\filldraw[fill=primary!14, draw=muted!35] (6,2) rectangle (7,3);
\node[text=black, font=\scriptsize] at (6.5,2.5) {1};
\filldraw[fill=primary!18, draw=muted!35] (7,2) rectangle (8,3);
\node[text=black, font=\scriptsize] at (7.5,2.5) {2};
\node[anchor=east, font=\scriptsize] at (-0.1,3.5) {RQ4 several surfaces};
\filldraw[fill=primary!18, draw=muted!35] (0,3) rectangle (1,4);
\node[text=black, font=\scriptsize] at (0.5,3.5) {2};
\filldraw[fill=primary!22, draw=muted!35] (1,3) rectangle (2,4);
\node[text=black, font=\scriptsize] at (1.5,3.5) {3};
\filldraw[fill=primary!22, draw=muted!35] (2,3) rectangle (3,4);
\node[text=black, font=\scriptsize] at (2.5,3.5) {3};
\filldraw[fill=white, draw=muted!35] (3,3) rectangle (4,4);
\filldraw[fill=primary!14, draw=muted!35] (4,3) rectangle (5,4);
\node[text=black, font=\scriptsize] at (4.5,3.5) {1};
\filldraw[fill=white, draw=muted!35] (5,3) rectangle (6,4);
\filldraw[fill=primary!14, draw=muted!35] (6,3) rectangle (7,4);
\node[text=black, font=\scriptsize] at (6.5,3.5) {1};
\filldraw[fill=primary!18, draw=muted!35] (7,3) rectangle (8,4);
\node[text=black, font=\scriptsize] at (7.5,3.5) {2};
\node[anchor=east, font=\scriptsize] at (-0.1,4.5) {Method precedent};
\filldraw[fill=primary!22, draw=muted!35] (0,4) rectangle (1,5);
\node[text=black, font=\scriptsize] at (0.5,4.5) {3};
\filldraw[fill=primary!22, draw=muted!35] (1,4) rectangle (2,5);
\node[text=black, font=\scriptsize] at (1.5,4.5) {3};
\filldraw[fill=primary!35, draw=muted!35] (2,4) rectangle (3,5);
\node[text=black, font=\scriptsize] at (2.5,4.5) {6};
\filldraw[fill=white, draw=muted!35] (3,4) rectangle (4,5);
\filldraw[fill=primary!14, draw=muted!35] (4,4) rectangle (5,5);
\node[text=black, font=\scriptsize] at (4.5,4.5) {1};
\filldraw[fill=white, draw=muted!35] (5,4) rectangle (6,5);
\filldraw[fill=primary!30, draw=muted!35] (6,4) rectangle (7,5);
\node[text=black, font=\scriptsize] at (6.5,4.5) {5};
\filldraw[fill=primary!85, draw=muted!35] (7,4) rectangle (8,5);
\node[text=white, font=\scriptsize] at (7.5,4.5) {18};
\node[anchor=north, font=\footnotesize] at (4.0,5.2) {App domain};
\node[anchor=south, rotate=90, font=\footnotesize] at (-4.21,2.5) {Part of the thesis RQ};
\end{tikzpicture}

\caption[Evidence map: parts of the RQ and app domain]{Evidence map of the
review evidence (n~=~47): number of studies per part of the thesis RQ and app
domain, coded by hand from the full text; a study can fall in several cells.}
\Description{Heat map with five rows (RQ1 to RQ4 and method precedent) and eight
domains. RQ1 and method precedent are filled across most domains; the RQ2 row
has entries only for studies outside health and in reproductive health.}
\label{fig:lit-map-rq}
\end{figure}

\paragraph{LQ1: harms and capabilities.} The review evidence documents
capabilities rather than harms as such: what data an adversary or a recipient
can obtain or infer, and from what. Table~\ref{tab:lit-capabilities} collects
them as the input to the harm trees. Location is inferable from shared activity
routes even when users hide their homes behind the privacy zones that fitness
platforms offer~\parencite{hassan2018analysis,dhondt2022run}; personality traits
are inferable from tracker data~\parencite{zufferey2023watch}; and health and
fitness apps pass data to third parties, often without saying so in their
policies~\parencite{huckvale2015unaddressed,huckvale2019assessment,grundy2019data,tangari2021mobile}.
For reproductive health, studies after the overturn of Roe v.~Wade add law
enforcement as a recipient: users were concerned about access by law
enforcement and third parties~\parencite{cao2024deleted}, and the policies of
female health apps named the possible disclosure of personal data to law
enforcement~\parencite{malki2024exploring}. Mental health apps showed
unnecessary permissions, insecure cryptography and leaks of personal data and
credentials~\parencite{iwaya2023privacy}. The survey of wearable activity
trackers~\parencite{niksirat2024wearable} systematises such risks for
trackers in general. The 154 LQ1 records still awaiting their full text
include many inference attacks from wrist-worn sensors (keystrokes, PINs,
handwriting) and studies of re-identification from activity and heart data;
once assessed, they will extend the table.

\begin{table}[tbp]
\centering
\begin{singlespace}
\caption[Capabilities documented in the review evidence]{Capabilities
documented in the review evidence (LQ1, full text assessed), with the data they
need and the actor who can use them. Input to the reference harm trees.}
\label{tab:lit-capabilities}
\small
\begin{tabularx}{\textwidth}{@{}>{\raggedright\arraybackslash}X>{\raggedright\arraybackslash}p{3.0cm}>{\raggedright\arraybackslash}p{2.8cm}>{\raggedright\arraybackslash}p{2.4cm}@{}}
\toprule
\textbf{Capability} & \textbf{Data needed} & \textbf{Actor} & \textbf{Studies} \\
\midrule
Locate sensitive places despite privacy zones & Shared activity routes & Other platform users & \cite{hassan2018analysis,dhondt2022run} \\
Infer personality traits & Tracker activity data & Recipients of the data & \cite{zufferey2023watch} \\
Profile users across apps & Identifiers and health data in app traffic & Third parties (advertising, analytics) & \cite{huckvale2015unaddressed,huckvale2019assessment,grundy2019data,tangari2021mobile} \\
Obtain reproductive data for prosecution & Cycle and pregnancy records & Law enforcement & \cite{cao2024deleted,malki2024exploring} \\
Read leaked personal data and credentials & Logs, web requests & Attackers & \cite{iwaya2023privacy} \\
\bottomrule
\end{tabularx}
\end{singlespace}
\end{table}

\paragraph{LQ2: communication.} The 21 LQ2 studies of the review evidence
analyse policies above all: their availability and
quality~\parencite{sunyaev2015availability,robillard2019availability,oloughlin2019reviewing},
their content for menstrual and female health apps~\parencite{shipp2020period,malki2024exploring}
and their match with app behaviour~\parencite{huckvale2019assessment,tangari2021mobile}.
Few go beyond the policy: Hutton et al.\ evaluated apps and policies against
heuristics~\parencite{hutton2018assessing}, Malki et al.\ combined policies
with an inspection of the apps~\parencite{malki2024exploring} and Hudig and
Singh studied the data flows of fertility trackers and users'
preferences~\parencite{hudig2025intimate}. One study rated apps' potential
damage from the information in their store listings; none analyses marketing,
and none codes which harms the communication names. Among the records awaiting
their full text, a discourse analysis of how period- and fertility-tracking
companies describe privacy and an analysis of store description and policy of
mental health apps come closest to this thesis; they must be assessed before
the gap can be stated as confirmed.

\paragraph{LQ3: language models.} The 14 LQ3 studies use language models to
annotate or assess privacy policies, for data practices, contextual
integrity, compliance with privacy law or GDPR transparency, and contradictions
within policies or with labels~\parencite{rodriguez2024llm,chanenson2025automating,xie2025evaluating}.
Each reports an empirical evaluation of the models, as LQ3 required; none of
them annotates harms or stance, which is the coding task of this thesis. Their validation designs are the precedents for the
LLM-assisted coding; the cautions of Baumann et al.\ and Törnberg on
configuration choices and validation remain in the contextual
collection~\parencite{baumann2025llm,tornberg2024best}.

\paragraph{LQ4: framing.} Six studies analyse how texts present risk or
protection: vagueness in policies~\parencite{lebanoff2018automatic}, hedging and
the downplaying of risk in breach notifications~\parencite{zou2019might}, the
transparency of genetic-testing policies~\parencite{doan2024transparency}, the
promises made for privacy-enhancing techniques~\parencite{martin2025cookies},
risk communication in intimate health technologies~\parencite{moniz2025toward}
and privacy washing~\parencite{brackin2026privacy}. Two concern health
(genetic testing and intimate health technologies), and none analyses the texts
of health tracking apps.

\paragraph{Answer to the review question and gaps.} The confirmed evidence
documents capabilities and harms for health tracking apps well, especially
location and attribute inference and third-party sharing, and it shows how
policies of such apps have been sampled and coded. It does not show how
providers word the harms they acknowledge, whether they downplay them, or how
they frame the protection of PETs: no study of the review evidence combines
harm coverage, stance and protection framing across several surfaces of health
tracking apps. The language-model studies validate data-practice annotation,
not the coding of harms. Three limits qualify this answer. First, the 47
studies come disproportionately from the Methods Digest (36), because only
their full texts were at hand; the 244 records awaiting their full text are the
larger and more recent part of the corpus. Second, the domain and surface codes
of the evidence map come from keyword rules on titles and abstracts and are
coarse. Third, the ACM Digital Library, Scopus, Web of Science and Google
Scholar were not searched, and five venues were filtered on titles only.
\section{Sample}
\label{sec:results-sample}
This section reports how the sample of Section~\ref{sec:app-sampling} was
drawn. All figures come from the collector outputs and the ranking script
(Appendix~\ref{app:scripts}); the sample is listed in
Appendix~\ref{app:app-sample}.
\todo{The screening decisions in 04\_Study/\allowbreak App-Sample/\allowbreak screening.csv were
prepared by an AI assistant from the store listings. Check every decision
before this section states who screened.}

\paragraph{Collection.} The App Store was searched on 23 October 2026 between
19:46 and 19:53~UTC\@. Of the 49 queries, 48 returned between 64 and 190 results;
\emph{reproduktive Gesundheit} returned none. The first 30 positions held 1{,}440
hits and 848 distinct listings. Google Play was searched on 24 October 2026
between 15:07 and 15:29~UTC\@. Each query returned between 8 and 50 results, 1{,}178
in total; the first 30 positions held 1{,}138 hits and 645 distinct apps, and
every hit carried an app ID\@. The product pages of all 645 apps were read on
24 and 25 October 2026; one page that timed out was read again, so none is
missing.

\paragraph{Preinstalled apps.} The vendor shares for September 2026, read on
22 October 2026, give Samsung 39.3\,\%, Apple 27.6\,\% and Xiaomi 21.4\,\% of
German mobile page views, and Huawei and Google about 1.5\,\% each
\parencite{statcounter2026vendor}. Apple is the only vendor of iOS (27.7\,\% of
page views). On Android (72.3\,\%), Samsung leads with 54\,\% of the
platform's page views, followed by Xiaomi with 30\,\%
\parencite{statcounter2026os}. The leading vendors are therefore Apple and
Samsung. Their qualifying apps are Apple Health and Apple Fitness on the iPhone
and Samsung Health on Samsung phones. The rule therefore yields three
preinstalled apps, which leaves 47 places for the ranking. A threshold of 10\,\%
of all page views would admit Xiaomi as well, but no preinstalled tracker could
be identified on Xiaomi's international software, so the three apps are the same
under that rule and under a rule that covers 80\,\% of Android. All three have store listings, which
were removed from the ranking; Apple Health, for instance, has 2{,}134 ratings
in the German App Store.

\paragraph{German-market popularity score.} Of the candidates, 125 apps were
found in both stores. The calibration used the 116 with positive rating and
review counts in both stores. Their counts correlated at Spearman 0.79, and
the median ratio was $m = 0.233815$. With the German platform shares of
September 2026, Android 72.3\,\% and iOS
27.7\,\%~\parencite{statcounter2026os}, the conversion factor was
\[
  c = 0.233815\times\frac{0.277}{0.723}\approx 0.089581.
\]
Thus, one App Store rating contributed approximately 0.09 score units and one
Play review contributed one unit. For example, komoot had 417{,}443 German App
Store ratings and 99{,}884 German Play reviews. Using the unrounded factor,
its score was $P_{\mathrm{komoot}}\approx137{,}279$. This is a combined feedback
score, not a user count; the script ranked candidates with unrounded scores.
Across all 644 Google Play candidates, the German review count and the
worldwide rating count correlated at 0.95, yet they rank the leading apps
differently: among the 100 apps with the most installs worldwide, the German
review count ranged from 5 to 99{,}884.

\paragraph{Screening and selection.} Table~\ref{tab:sample-flow} gives the
cascade. Merging the two stores' listings gave 1{,}365 candidates. In descending
order of German-market popularity score, 70 candidates were reached before
the 47 places were filled.
Of these, 11 were excluded at screening because the app records no health or
body data over time (E1), and 12 were skipped because their developer was
already in the sample. Developer names that differ between the two stores were
matched by hand for 13 providers, among them Clue, Pregnancy+, Fddb and
Headspace, each on the evidence listed in Table~\ref{tab:developer-aliases};
no store account was found under two providers, and no two providers among the
70 candidates reached named the same web domain. Three apps whose titles were
translated between the stores were paired by hand. Across the main selection and all sensitivity checks, 275
listings were screened: 156 were eligible and 119 excluded (E1~114, E2~1,
E3~2, E4~1, E5~1).

\begin{table}[t]
\centering
\begin{singlespace}
\caption[From search hits to the 47 ranked apps]{From search hits to the 47 ranked apps. Positions are the first 30 per
query; the same app in both stores is one candidate.}
\label{tab:sample-flow}
\small
\begin{tabular}{@{}lrr@{}}
\toprule
 & \textbf{App Store} & \textbf{Google Play} \\
\midrule
Queries & 49 & 49 \\
Hits within the first 30 positions & 1{,}440 & 1{,}138 \\
Distinct listings & 848 & 645 \\
Listings of preinstalled apps removed & 2 & 1 \\
Listings ranked & 846 & 644 \\
\midrule
\multicolumn{3}{@{}l}{\textbf{Candidates after merging the stores}} \\
\quad in both stores & \multicolumn{2}{r}{125} \\
\quad App Store only & \multicolumn{2}{r}{721} \\
\quad Google Play only & \multicolumn{2}{r}{519} \\
Reached in descending score order & \multicolumn{2}{r}{70} \\
\quad excluded at screening (E1) & \multicolumn{2}{r}{11} \\
\quad skipped: developer already in the sample & \multicolumn{2}{r}{12} \\
\quad selected & \multicolumn{2}{r}{47} \\
\bottomrule
\end{tabular}
\end{singlespace}
\end{table}

\paragraph{The sample.} The 47 ranked apps come from 47 providers; 33 are listed
in both stores, 8 only on Google Play and 6 only on the App Store. With Apple
Health, Apple Fitness and Samsung Health, the sample holds 50 apps. By primary
category, the ranked apps comprise 22 activity apps, 12 diet and weight apps, 4
cycle trackers, 2 pregnancy apps, 3 sleep apps, 2 mood apps, 1 heart-rate app and
1 wearable companion. The selected apps have between 648 and 447{,}096 German App
Store ratings (39 apps with an App Store listing) and between 1{,}486 and
99{,}884 German Google Play reviews (41 apps with a Google Play listing). All are
free to install, and all 41 Google Play listings name a privacy policy. Of the
47 apps, 42 were found within the first ten positions of at least one query and
5 within positions 11 to 20.

\paragraph{Sensitivity.} Table~\ref{tab:sensitivity} shows how many of the 47
selected providers change when one rule is changed at a time. Without the
developer rule, Leap Fitness Group would take 7 of the 47 places, Simple Design
3, and Google and Nike 2 each (Table~\ref{tab:developer-rule}).

\begin{table}[t]
\centering
\begin{singlespace}
\caption[Sensitivity of the selection]{Sensitivity of the selection. Number of the 47 selected providers that
are not selected when one rule changes.}
\label{tab:sensitivity}
\small
\begin{tabularx}{\textwidth}{@{}Xr@{}}
\toprule
\textbf{Variant} & \textbf{Providers changed} \\
\midrule
No developer rule (the 47 places go to 37 providers) & 10 \\
Equal weights for iOS and Android & 5 \\
Alternative scale (share of each store's total count) & 0 \\
Worldwide Google Play rating count instead of German reviews & 11 \\
Earlier selection: store alternation with the worldwide Google Play count & 10 \\
Depth 20 instead of 30 & 1 \\
App Store searched to depth 200 & 2 \\
English keywords only & 2 \\
German keywords only & 3 \\
One keyword left out (49 runs; largest change) & 1 \\
\bottomrule
\end{tabularx}
\end{singlespace}
\end{table}

\section{What Privacy Communication Covers}
\todo[inline]{To be written.}

\section{What Privacy Communication Omits}
\todo[inline]{Inference, breaches, incident handling.}

\section{Differences Across Platforms and Categories}
\todo[inline]{To be written.}

\section{Privacy Content in Advertising}
\todo[inline]{To be written.}
